\section{Round-sixteen positive closure: adapted latent phases and exact nonlinear aggregation}
\label{sec:r16-d1}

A phase is not a terminal event read from the complete future path.  It is a
positive latent variable chosen at time zero on the microscopic probability
space.  Its posterior is part of the adapted state.  Component laws, their
large deviations, response charts, and Gaussian tangents are derived from the
joint labelled law.  The restricted-pressure sign is fixed throughout.

\subsection{Positive latent phase labels}

For each finite system let \(J_\varepsilon\in\{0,1,\ldots,m\}\) be a positive
measurable initial label.  The events are defined by a partition of unity on an
initial slow order parameter: labels \(1,\ldots,m\) select isolated regular
basins and label zero is the transition/boundary component.  Let
\[
 w_{\varepsilon,j}=\mathbb P(J_\varepsilon=j),
 \qquad
 \alpha_{\varepsilon,j}=-{1\over\mu_\varepsilon}
 \log w_{\varepsilon,j},
\]
and let \(P_{\varepsilon,j}\) be the conditional path law.  The label is
\(\mathcal F_0\)-measurable and therefore unchanged by restriction to any
later horizon.  Its posterior
\[
 \pi_t(j)=\mathbb P(J_\varepsilon=j\mid\mathcal F_t^{\rm obs})
\]
is adapted and obeys Bayes' rule.

\begin{theorem}[Horizon-consistent positive phase process]
\label{thm:r16-d1-phase}
The joint laws of \((J_\varepsilon,X^\varepsilon_{[0,T]})\) are projectively
consistent in \(T\).  Conditional transition kernels are causal on the
augmented state \((\pi_t,X_t^{\rm resolved})\), and
\[
 \pi_t(j)={w_{\varepsilon,j}L_{t,j}
 \over\sum_k w_{\varepsilon,k}L_{t,k}},
\]
where \(L_{t,j}\) is the exact complete-filtration likelihood projected to the
observation filtration.  No full-horizon order parameter is used to label an
intermediate state.
\end{theorem}

\begin{proof}
The label is sampled from the initial sigma-field and the platform dynamics is
then run conditionally on it.  Restricting a path does not change its initial
label, giving projective consistency.  Disintegration of the joint law gives
the component transition kernels.  Bayes' formula gives the posterior; its
normalizing denominator is the positive mixture likelihood.  Including the
posterior in the state makes the transition adapted and Markov at the law
level.
\end{proof}

\subsection{Correct component pressures}

For a path source \(F\), define
\[
 \widetilde Q_{\varepsilon,j}(F)
 ={1\over\mu_\varepsilon}
 \log\mathbb E[e^{\mu_\varepsilon F}\mid J_\varepsilon=j]
\]
and the unnormalized component pressure
\[
 Q^{\rm un}_{\varepsilon,j}(F)
 ={1\over\mu_\varepsilon}
 \log\mathbb E[e^{\mu_\varepsilon F};J_\varepsilon=j].
\]

\begin{lemma}[Restricted-pressure sign]
\label{lem:r16-d1-sign}
Exactly at finite volume,
\[
 Q^{\rm un}_{\varepsilon,j}(F)
 =-\alpha_{\varepsilon,j}
  +\widetilde Q_{\varepsilon,j}(F).
\]
Consequently
\[
 Q_\varepsilon(F)
 ={1\over\mu_\varepsilon}
 \log\sum_j
 \exp\{\mu_\varepsilon[-\alpha_{\varepsilon,j}
 +\widetilde Q_{\varepsilon,j}(F)]\}.
\]
The phase cost occurs once.
\end{lemma}

\begin{proof}
Factor the expectation as
\(w_{\varepsilon,j}\mathbb E_j e^{\mu F}\) and take its logarithm.  Since
\(\mu^{-1}\log w_j=-\alpha_j\), the sign follows.  Summing the positive
component expectations gives the second identity.
\end{proof}

\subsection{Joint labelled LDP}

Assume the platform-specific A3 or B2 path law has a good rate and the initial
partition functions have limits
\(\alpha_j=\lim\alpha_{\varepsilon,j}\).  The phase partition is chosen inside
regular basin interiors with a boundary label retaining the transition band.

\begin{theorem}[Joint phase--path LDP]
\label{thm:r16-d1-joint}
The joint variable \((J_\varepsilon,X^\varepsilon)\) satisfies a good LDP on
the discrete-label product with rate
\[
 \mathcal I(j,x)=\alpha_j+I_j(x),
\]
where \(I_j\) is the good rate of the genuine conditional law
\(P_{\varepsilon,j}\).  Each component lower bound is proved by recovery
inside its positive initial basin; the boundary label has its own conditional
rate.  Contraction gives the physical rate
\[
 I(x)=\min_j\{\alpha_j+I_j(x)\}.
\]
\end{theorem}

\begin{proof}
For a closed set, split the joint probability over the finitely many positive
labels and use the conditional upper bounds together with
\(w_{\varepsilon,j}=e^{-\mu_\varepsilon\alpha_{\varepsilon,j}}\).  For an open labelled set, choose one component and
multiply its conditional recovery probability by its label weight.  The
conditional lower recovery is obtained by first choosing the initial point in
the interior of the corresponding basin and then applying the A3 or B2
platform recovery; openness prevents escape to another label.  The boundary
band is retained rather than assigned to a neighboring phase.  Finite-label
exponential tightness gives goodness.  Forgetting the label gives the minimum
by ordinary contraction.
\end{proof}

\subsection{Component analytic charts}

A real rate gap is not used as a complex zero-free theorem.  Each component has
its own positive restricted transfer operator or exact-number coefficient.

\begin{theorem}[Phase-local zero-free response charts]
\label{thm:r16-d1-charts}
On every regular component there is a complex source neighborhood on which the
normalized component partition function is zero free and
\(\widetilde Q_{\varepsilon,j}\) converges with all fixed derivatives.  On the
Sinai platform this follows from the simple eigenvalue of the component killed
return operator; on the hard-sphere platform it follows from the B1
exact-number coefficient theorem with the positive initial basin insertion.
The boundary component is retained but no simple analytic chart is asserted
at a coexistence singularity.
\end{theorem}

\begin{proof}
For Sinai, positivity and a spectral gap of the basin-restricted returned
operator give a simple leading eigenvalue separated by a uniform contour.
Analytic perturbation and Rouch\'e's theorem give a zero-free component
partition function.  For hard spheres, insert the positive basin cutoff in the
exact-\(N\) canonical law.  B1's angular gap and continuous Fourier theorem
hold uniformly on compact basin interiors, giving a nonvanishing saddle
coefficient and derivative convergence.  These are quantitative complex
estimates for component objects, not consequences of a real LDP gap.
\end{proof}

\subsection{Nonlinear phase semigroup}

Let \(S_{t,j}\) be the proved component semigroup on the platform state.  On
the augmented posterior state define
\[
 (\mathcal S_t\Phi)(\pi,x)
 ={1\over\theta}
 \log\sum_j\pi(j)
 \exp\{\theta S_{t,j}\Phi_j(x)\}.
\]

\begin{theorem}[Adapted log-sum-exp tower]
\label{thm:r16-d1-semigroup}
The phase-posterior value satisfies dynamic programming on the augmented state.
At finite volume it equals the exact mixture log-Laplace value, and its
large-speed limit is
\[
 \max_j\{-\alpha_j+V_j\}.
\]
At ties the cotangent is the convex hull of the active labelled phase laws,
and posterior weights are determined by the retained subexponential component
normalizers.
\end{theorem}

\begin{proof}
Condition first on the initial latent label and apply the component tower.
Bayes' formula updates \(\pi_t\); substituting it in the next time block
cancels the intermediate evidence denominator and yields the exact augmented
tower.  The finite-volume identity is Lemma~\ref{lem:r16-d1-sign} applied to
the component values.  Laplace's principle gives the maximum.  Differentiating
finite log-sum-exp gives the softmax mixture of component tangents; along a tie,
subsequential limits retain the component prefactors and give the active convex
hull.
\end{proof}

\subsection{Labelled Gaussian tangents}

\begin{theorem}[Gaussian mixtures at coexistence]
\label{thm:r16-d1-gaussian}
On each regular phase chart, the centered component field converges jointly
with its likelihood ratio to the platform Gaussian tangent with mean and
covariance \((m_j,\Sigma_j)\).  Jointly with the latent label,
\[
 (J_\varepsilon,Z_\varepsilon)
 \Longrightarrow\sum_j\omega_j\,
 \delta_j\otimes N(m_j,\Sigma_j),
\]
where \(\omega_j\) is the limit of the posterior/component prefactors along
the chosen sequence.  At a tie there is no single covariance or single
Cameron--Martin space; the labelled mixture retains all component geometries.
\end{theorem}

\begin{proof}
The A4 or B3 component central limit theorem and the zero-free chart give
conditional Gaussian convergence.  Multiplication by the positive label
weights gives joint convergence on the finite discrete product.  Every
subsequence has a further subsequence on which the normalized component
prefactors converge to \(\omega_j\).  Summing the conditional characteristic
functions gives the displayed labelled mixture.  The componentwise
Cameron--Martin likelihoods remain attached to their labels.
\end{proof}

\subsection{Conditioning and calibration}

\begin{theorem}[Phase-aware contractions]
\label{thm:r16-d1-contraction}
Source-dependent microcanonical conditioning, continuous observable
contraction, pressure differentiation, Gaussian tangent formation, and the
adapted log-sum-exp semigroup commute on the joint labelled law.  The initial
constraint saddle is solved within each positive component before nonlinear
aggregation.  A mechanical preparation field equals an entropic valuation
coefficient only under the explicit canonical-cocycle compatibility axiom.
\end{theorem}

\begin{proof}
Each operation is performed first on the genuine conditional component law.
B1 gives the component source-dependent saddle, the contraction theorem gives
component rates, and the component pressure Hessian gives its Gaussian
covariance.  Lemma~\ref{lem:r16-d1-sign} then aggregates the positive
unnormalized values with the phase cost counted once.  Since the label is part
of the state, all time operations are adapted.  Finally, equality of a
preparation multiplier and a valuation coefficient follows only by equality
of their likelihood cocycles on a nonconstant work observable; phase
decomposition alone supplies no such equality.
\end{proof}

\begin{theorem}[Round-sixteen D1 closure]
\label{thm:r16-d1-main}
The standalone phase theorem now uses an initial positive latent label,
horizon-consistent posterior dynamics, the correct restricted-pressure sign,
genuine conditional component LDPs and zero-free charts, exact nonlinear
log-sum-exp aggregation, and a labelled Gaussian mixture at coexistence.
\end{theorem}

\begin{proof}
Combine Theorems~\ref{thm:r16-d1-phase},
\ref{thm:r16-d1-joint}, \ref{thm:r16-d1-charts},
\ref{thm:r16-d1-semigroup}, \ref{thm:r16-d1-gaussian}, and
\ref{thm:r16-d1-contraction}, and
Lemma~\ref{lem:r16-d1-sign}.
\end{proof}
